\section{Introducción: qué es el Post-Training}

Esta clase la imparten investigadores de Liquid AI (autores de «LLM Engineer's Handbook»), quienes presentan de manera sistemática el flujo de Post-Training (posentrenamiento) de los modelos de lenguaje de gran escala (LLM). El Post-Training abarca todos los pasos de entrenamiento posteriores al preentrenamiento, y su objetivo es convertir un base model que solo sabe predecir el siguiente token en un AI assistant realmente útil.

Durante la etapa de preentrenamiento, el modelo se entrena sobre una cantidad masiva de texto original con el objetivo de next-token prediction; el base model que resulta tiene una gran capacidad de modelado del lenguaje, pero carece de la capacidad de responder preguntas o seguir instrucciones. El Post-Training es precisamente el paso clave que cierra esta brecha.

\begin{importantbox}{Las dos etapas centrales del Post-Training}
\begin{enumerate}
    \item \textbf{Supervised Fine-Tuning (SFT)}: mediante pares de instrucción y respuesta (instruction-answer pairs) se le enseña al modelo a seguir instrucciones y a responder preguntas en forma de conversación. En esta etapa el modelo aprende la estructura del chat template (system / user / assistant).
    \item \textbf{Preference Alignment (alineación de preferencias)}: se ajusta el comportamiento del modelo comparando la chosen answer (respuesta deseada) con la rejected answer (respuesta no deseada), de modo que las respuestas producidas se ajusten mejor a las preferencias humanas.
\end{enumerate}
\end{importantbox}

\begin{knowledgebox}{Diferencia de escala de datos entre Pre-Training y Post-Training}
El preentrenamiento utiliza decenas de billones (trillions) de tokens de texto original, con un conjunto de datos de tamaño extremadamente grande. En cambio, el conjunto de datos del Post-Training es mucho más pequeño, por lo general del orden de miles a millones de ejemplos, pero exige una \textbf{calidad} de los datos mucho mayor que la cantidad. Esta diferencia es la base para entender la metodología del Post-Training.
\end{knowledgebox}

\subsection{Resumen de la sección}

El Post-Training es el paso central que convierte un base model en un assistant utilizable, e incluye dos etapas principales: SFT y Preference Alignment. Se caracteriza por tener un conjunto de datos mucho más pequeño que el del preentrenamiento, pero con exigencias de calidad muy altas.

\newpage
